\documentclass[10pt,a4paper]{amsart}
\usepackage[margin=19mm]{geometry}
\usepackage{amsmath,amssymb,mathtools}
\usepackage[scr=rsfs,cal=euler]{mathalfa}
\usepackage{xfrac}
\usepackage[unicode,hidelinks,bookmarks=false]{hyperref}
\newcommand{\ie}{i.e.\ }
\newcommand{\cf}{cf.\ }
\newcommand{\resp}{respectively\ }
\newcommand{\Subsection}{\subsection}
\providecommand{\nobreakdash}{\nobreak\hbox{-}}
\setlength{\emergencystretch}{2em}
\multlinegap0pt
\usepackage{amssymb}
\newtheorem{lemma}{Lemma}
\title{Weak central polynomials for algebras of multiplications of simple algebras}
\author{Vesselin Drensky, Mikhail Zaicev}
\date{Source: arXiv:2609.10362v1 (CC BY 4.0)}
\begin{document}
\maketitle

\noindent\emph{Source and licence.} Weak central polynomials for algebras of multiplications of simple algebras. Version arXiv:2609.10362v1 (CC BY 4.0), distributed by arXiv under the Creative Commons Attribution 4.0 International licence, http://creativecommons.org/licenses/by/4.0/, as recorded in the arXiv licence metadata for this exact version and shown on the abstract page https://arxiv.org/abs/2609.10362v1. Full licence notice: \texttt{licenses/arxiv-2609.10362-CC-BY-4.0.txt}.\par\smallskip

\noindent\textit{Editorial scope.} Counted unit: the article's lemma lemma (source0.tex lines 122--125). The article's complete original proof of it is delivered with it (source0.tex lines 127--148, one \texttt{proof} environment of 22 lines). No mathematics is written, repaired or re-derived: the statement and the proof are the article's own lines, in the article's own notation, and no retained line is substituted or re-flowed.  No further source-local context is needed: every cross-reference of every retained line resolves inside the two delivered spans.  Nothing was omitted from any retained span, so the package carries no omission record.  No retained line contains a citation, so the wrapper carries no bibliography and no bibliography entry.  No retained line contains a figure environment, an image-inclusion command or drawing code, so no image is delivered and the package declares no asset dependency.  Editorial formatting only, in addition: an AMS \texttt{amsart} preamble that restates the article's own declarations and macros used by the retained lines, namely the environments \texttt{lemma} on separate counters, together with the standard packages those lines need.  The retained environments print in this document as Lemma 1 in order of appearance, whereas the article prints the same items with the numbering of its own full document.  The complete native source of the article as distributed in the arXiv e-print for this exact version is bundled unchanged as \texttt{sources/209904/source0.tex}.
\begin{lemma}\label{the main lemma}
Let $c(X)=c(x_1,\ldots,x_n)\in F\langle X\rangle$ be a non-trivial multilinear central polynomial of the algebra $R$.
Then there exists a polynomial $c'(Y)\in F\langle Y\rangle$ which is a non-trivial central polynomial for the pair $(R,V)$.
\end{lemma}

\begin{proof}
Let $V$ have a vector space basis $U=\{u_1,u_2,\ldots\}$. Since $V$ generates $R$ as an algebra,
$R$ is spanned on the products
\[
W=\{w_i=u_{i_1}\cdots u_{i_h}\mid u_{i_j}\in U,h=1,2,\ldots\}.
\]
Let us fix a subset $W_0$ of $W$ which is a basis of the vector space $R$.
Since $c(X)$ is not a polynomial identity for $R$ and is multilinear, there are elements $w^{(1)},\ldots,w^{(n)}$ from $W_0$ such that
$c(w^{(1)},\ldots,w^{(n)})\not=0$. Let
\[
w^{(p)}=u_{i_1}^{(p)}\cdots u_{i_{h_p}}^{(p)},\quad u_{i_j}^{(p)}\in U.
\]
We define
\[
c'(Y)=c(y_{11}\cdots y_{1h_1},\ldots,y_{n1}\cdots y_{nh_n})\in F\langle Y\rangle.
\]
Evaluated on $V$, the products $y_{j1}\cdots y_{jh_j}$ belong to $R$. Hence $c'(Y)$ is a central element evaluated on $V$, i.e., $c'(Y)$
is a central polynomial for the pair $(R,V)$. It is non-trivial because
\[
c(u_{i_1}^{(1)}\cdots u_{i_{h_1}}^{(1)},\ldots,u_{i_1}^{(n)}\cdots u_{i_{h_n}}^{(n)})=c(w^{(1)},\ldots,w^{(n)})\not=0.
\]
\end{proof}

\end{document}
